%%%%%%%%%%%%%%%%%%%%%%%%%%%%%%%%%%%%%%%%%%%%%%%%%%%%%%%%%%%%%%
\chapter{Rewritten Nets and Control-Flow Replica}
\label{app:replica}
\label{app:rewritten}
%%%%%%%%%%%%%%%%%%%%%%%%%%%%%%%%%%%%%%%%%%%%%%%%%%%%%%%%%%%%%%
% Sources: trucksimulation/process_discovery/outputs/runs/e2_split13apr_20261106 (repair, rewritten nets);
%   simulation_e3/runs/split13apr_20261106/{main,noguard,guard_orig,samelang,ablation_flower},
%   runs/hm07_raw_20261111/{noguard,guard,sm_noguard}; tables: process_discovery/make_rq1_tables.py (app_rewritten_*);
%   replica: process_discovery/ablation_silent_guard.py --reference train -> analysis/ablation_silent_guard_20261113,
%   table: process_discovery/make_replica_table.py -> tables/app_control_flow_replica.tex

In an earlier stage of this work, the admissible nets were not simulated as
discovered. They were repaired and rewritten first, and the stop guard acted only
on visible transitions. This appendix reports that stage, because it explains why
the final design looks ahead on the discovered net (Section~\ref{sec:exp_guard})
and what a simulator does with the language of a net.

\section{Repair and Rewriting}
\label{app:rewritten_repair}
The languages of IM 0.0 and ILP 1.00 were restricted exactly to rules 3 and 5:
the reachability graph of each net was combined with automata for the rules, and
states with the same continuations were merged. Each of the three admissible nets
was then rewritten without silent transitions, with one place for each state of
the visit and one transition for each activity that can happen in that state. The
language stayed the same; an activity therefore appears on several transitions.
The repair kept every observed sequence and raised the precision of IM 0.0 from
0.842 to 0.888 and that of ILP 1.00 from 0.773 to 0.825
(Table~\ref{tab:app_rewritten_funnel}).

\begin{table}[htbp]
\centering
\footnotesize
\setlength{\tabcolsep}{3pt}
\begin{tabular}{llrrrr}
\toprule
 & & \textbf{Sequences} & \multicolumn{2}{c}{\textbf{Precision (training)}} & \textbf{Gen.} \\
\textbf{Net} & \textbf{Rules violated} & \textbf{with 4 stops} & \textbf{discovered} & \textbf{state machine} & \textbf{(of 16)} \\
\midrule
HM 0.7 & 6 (bound) $\rightarrow$ 6 (bound) & 27 $\rightarrow$ 27 & 0.998 & 0.998 & 2 $\rightarrow$ 2 \\
IM 0.0 & 3, 5, 6 (bound) $\rightarrow$ 6 (bound) & 2009 $\rightarrow$ 1583 & 0.842 & 0.888 & 13 $\rightarrow$ 13 \\
ILP 1.00 & 3, 5, 6 (bound) $\rightarrow$ 6 (bound) & 6561 $\rightarrow$ 3937 & 0.773 & 0.825 & 16 $\rightarrow$ 16 \\
\bottomrule
\end{tabular}
\caption{The three nets that pass A1 and A2, as discovered and as state machines of the language restricted to rules 3 and 5 (control of Appendix~\ref{app:replica}). Rules violated: rules whose violation the language permits, discovered $\rightarrow$ restricted (6 (bound): more than four stops, enforced by the stop guard). Sequences with 4 stops: distinct sequences the language permits (5 occur in training). Precision: ETConformance in ProM 6.15. Gen.: new test sequences accepted, discovered $\rightarrow$ restricted.}
\label{tab:app_rewritten_funnel}
\end{table}


\paragraph{Precision and simulation.} Simulated with the same settings as in E3,
the three rewritten nets generated the same distribution of sequences (sequence
TVD 0.015 to the training log for all three), although their precision ranged
from 0.825 to 0.998 (Table~\ref{tab:app_rewritten_matrix}). The behaviour that
precision counts against a net was hardly generated, because the routing is
learned from the observed visits. The exception is the unrepaired language of
ILP 1.00: rewritten, it has a single state for all stops, from which every stop
and Gate Out can follow (precision 0.773). It combined the stops wrongly
(sequence TVD 0.077, 0.36\,\% unseen sequences, against 0.015 and 0.01 to
0.05\,\%).

\begin{table}[htbp]
\centering
\scriptsize
\setlength{\tabcolsep}{2.5pt}
\begin{tabular}{lrcllrr|lr}
\toprule
 & & & \multicolumn{4}{c|}{\textbf{Decision (against the training log)}} & \multicolumn{2}{c}{\textbf{Reported (test log)}} \\
\textbf{Net} & \textbf{Prec.} & \textbf{Gate} & \textbf{Sequence TVD} & \textbf{Stop-count TVD} & \textbf{Overr.} & \textbf{Unseen [\%]} & \textbf{Sequence TVD} & \textbf{TTT med./P90} \\
\midrule
HM 0.7 & 0.998 & yes & 0.015 [0.013, 0.018] & 0.004 [0.003, 0.007] & 117 & 0.01 & 0.096 [0.094, 0.097] & 58 / 136 \\
IM 0.0 & 0.888 & yes & 0.015 [0.014, 0.016] & 0.002 [0.001, 0.002] & 0 & 0.04 & 0.097 [0.093, 0.100] & 58 / 141 \\
ILP 1.00 & 0.825 & yes & 0.015 [0.013, 0.017] & 0.002 [0.001, 0.002] & 4 & 0.05 & 0.098 [0.096, 0.099] & 59 / 141 \\
\bottomrule
\end{tabular}
\par\smallskip{\footnotesize\raggedright Each net is repaired, rewritten without silent transitions and run with the stop guard on the visible transitions. Means over five seeds with [minimum, maximum]. Prec.\ is ETConformance on the rewritten net (training). Gate means no visit without a stop, no overlapping stops of one truck and at most four stops, in every seed. TVD is the total-variation distance to the training log, for the sequences and for the number of stops per visit. Overr.\ counts the guard overrides per run. Unseen are generated visits whose sequence does not occur in the training log. The right part compares with the test log, whose TTT median is 33 and P90 75 minutes.\par}
\caption{Comparison of the three admissible nets as rewritten state machines.}
\label{tab:app_rewritten_matrix}
\end{table}


\paragraph{The simple guard and silent decisions.} On the discovered HM 0.7, the
simple guard, which masks only the visible stops and Gate Out, changed no draw: the
decisions are taken by silent transitions, and when a stop or Gate Out becomes
enabled it is often the only step left (449 overrides per run). On the rewritten
net it changed 3,206 draws per run, but still had to give up 117 times, in states
from which the plan could no longer be completed
(Table~\ref{tab:app_rewritten_2x2_hm}). The sequences were the same in all four
conditions. The look-ahead guard of the final design reaches the same on the
discovered net without rewriting it, and never has to give up (E3).

\begin{table}[htbp]
\centering
\scriptsize
\setlength{\tabcolsep}{3pt}
\begin{tabularx}{\textwidth}{L{3.0cm}YYYYYYY}
\toprule
\textbf{HM 0.7 in ProSiT} & \textbf{Without stop [\%]} & \textbf{Overlap [\%]} & \textbf{5+ stops [\%]} & \textbf{Unseen [\%]} & \textbf{Seq.\ TVD} & \textbf{Stop-count TVD} & \textbf{Guard: changed / overr.} \\
\midrule
Original net, no guard & 0.00 & 0.00 & 0.017 & 0.080 & 0.016 & 0.003 & -- \\
Original net, guard & 0.00 & 0.00 & 0.011 & 0.054 & 0.016 & 0.005 & 0 / 449 \\
Rewritten net, no guard & 0.00 & 0.00 & 0.013 & 0.075 & 0.015 & 0.005 & -- \\
Rewritten net, guard & 0.00 & 0.00 & 0.000 & 0.010 & 0.015 & 0.004 & 3,206 / 117 \\
\bottomrule
\end{tabularx}
\caption{$2\times2$ contrast for HM 0.7 in E3 (Table~\ref{tab:exp_e3_contrast}): original net (22 places, 45 transitions, 34 silent) or the same language rewritten without silent transitions, stop guard on the visible transitions off or on; mean over five seeds, against the training log. Overlap: share of the multi-stop visits with overlapping stops. Unseen: generated visits whose sequence does not occur in the training log. Guard: draws whose weights the guard changed, and draws it could not decide (overrides), per run.}
\label{tab:app_rewritten_2x2_hm}
\end{table}


\paragraph{Parallel branches.} For IM 0.0, whose discovered net expresses the
free order of the stops by parallel branches that silent transitions can skip,
the representation decided much more (Table~\ref{tab:app_rewritten_2x2_im}; test
log). The discovered net generated 15.1\,\% visits without a stop and overlapping
stops in 22.7\,\% of the multi-stop visits; the same language rewritten without
silent transitions generated none.

\begin{table}[htbp]
\centering
\scriptsize
\setlength{\tabcolsep}{3pt}
\begin{tabularx}{\textwidth}{L{3.4cm}YYYYYY}
\toprule
\textbf{IM 0.0, simulated with ProSiT} & \textbf{Visits without stop [\%]} & \textbf{Overlapping stops [\%]} & \textbf{Visits with 5+ stops [\%]} & \textbf{Unseen sequences [\%]} & \textbf{Sequence TVD} & \textbf{Stop-count TVD} \\
\midrule
Observed (test log) & 0.00 & 0.00 & 0.000 & 0.00 & 0.000 & 0.000 \\
\midrule
Naive: original net, no guard & 15.07 & 22.71 & 0.031 & 19.02 & 0.360 & 0.277 \\
Guard only: original net + guard & 15.08 & 21.30 & 0.022 & 17.62 & 0.323 & 0.239 \\
Repair only: repaired net, no guard & 0.00 & 0.00 & 0.061 & 0.53 & 0.103 & 0.027 \\
Both: repaired net + guard & 0.00 & 0.00 & 0.000 & 0.05 & 0.097 & 0.021 \\
\midrule
Control: same language without silent transitions, no guard & 0.00 & 0.00 & 0.061 & 0.50 & 0.100 & 0.027 \\
\bottomrule
\end{tabularx}
\caption{Effect of the rule treatment on the simulated visits (E3), test period, means over five seeds. Original net: IM 0.0 as discovered (parallel branches that silent transitions can skip). Repaired net: language restricted to rules 3 and 5, rewritten without silent transitions. Guard: stop guard on the visible transitions. Control: the unrepaired language rewritten without silent transitions. Overlapping stops: share of the visits with more than one stop. Unseen sequences: generated visits whose sequence never occurs in the training log. TVD: total-variation distance to the distribution of the observed test log (0 = identical).}
\label{tab:app_rewritten_2x2_im}
\end{table}


\paragraph{E4 and E5 with two nets.} In this stage, E4 and E5 were run with the
rewritten HM 0.7 and with the rewritten IM 0.0, the net that a selection by
simulation measures would have chosen. Both needed the same four changes to the
queue and the service, the same configuration was chosen, the same acceptance
criterion was missed narrowly, and the scenario S1 shortened the moved visits by
9.3 minutes with both nets.

\paragraph{Outdated enabling times.} The rewritten IM 0.0 has seven loops, against
three for HM 0.7. In the earlier comparison, its final E4 configuration reached a
smaller cycle-time distance than that of HM 0.7 (6.3 against 7.9 minutes on the
test period). The advantage came from a defect of the training data: ProSiT keeps
the enabling time of a transition until it fires, also across loops, so the
waiting before Gate Out of multi-stop visits was measured from the first stop
(mean target 32.6 minutes, against 12.7 observed from the last stop). With the
enabling time reset whenever its token is consumed, the advantage disappeared
(cycle-time distance 8.5 against 8.0 minutes; training period 4.2 against 4.3),
while HM 0.7 changed by no more than the range of the seeds
(Section~\ref{sec:res_inspection}).

\section{Control-Flow Replica}
\label{app:replica_token_game}
To see why the representation matters, the routing of ProSiT was rebuilt as a
token game without time and resources. In every marking, one enabled transition
is drawn with a probability proportional to its weight, as in stochastic labelled
Petri nets \cite{burke2020stochasticPD,mannhardt2023stochasticProcessesModelling}.
The weights are the probability that a transition fires when it is enabled,
estimated by counting on the aligned training log. Every training visit is
simulated with its own manifest, with three seeds. The replica reproduces ProSiT
on the discovered IM 0.0: 15.7\,\% of its visits have no stop, against 15.1\,\% in
ProSiT. Table~\ref{tab:app_control_flow_replica} gives the results.

\begin{table}[htbp]
\centering
\scriptsize
\setlength{\tabcolsep}{2.5pt}
\begin{tabular}{lllrrrrr}
\toprule
\textbf{Net} & \textbf{Routing weights} & \textbf{Guard} & \textbf{No stop [\%]} & \textbf{Stops = plan [\%]} & \textbf{Unseen [\%]} & \textbf{Sequence TVD} & \textbf{Fallb./1000} \\
\midrule
\multicolumn{8}{l}{\textit{Weights and representation, no guard}} \\
IM 0.0, discovered & frequency & none & 26.8 & 42.3 & 33.769 & 0.489 & -- \\
IM 0.0, discovered & data + history & none & 15.7 & 76.1 & 18.163 & 0.232 & -- \\
IM 0.0, discovered & + marking & none & 0.0 & 99.8 & 0.002 & 0.004 & -- \\
IM 0.0, same language, rewritten & data + history & none & 0.0 & 99.8 & 0.011 & 0.005 & -- \\
IM 0.0, repaired and rewritten & data + history & none & 0.0 & 99.8 & 0.011 & 0.005 & -- \\
ILP 1.00, discovered & frequency & none & 47.8 & 24.0 & 61.228 & 0.708 & -- \\
Flower (= ILP 1.00 rewritten) & frequency & none & 48.0 & 23.9 & 61.328 & 0.709 & -- \\
Flower (= ILP 1.00 rewritten) & data + history & none & 0.0 & 99.8 & 0.012 & 0.005 & -- \\
HM 0.7, discovered & data + history & none & 0.0 & 99.9 & 0.000 & 0.010 & -- \\
HM 0.7, rewritten & data + history & none & 0.0 & 99.9 & 0.000 & 0.009 & -- \\
\midrule
\multicolumn{8}{l}{\textit{Guard variants, weights data + history}} \\
IM 0.0, discovered & data + history & local (thesis) & 15.7 & 77.1 & 17.501 & 0.225 & 237.7 \\
IM 0.0, discovered & data + history & look-ahead & 0.0 & 100.0 & 0.156 & 0.167 & 0.0 \\
IM 0.0, discovered & data + history & rejection & 0.0 & 100.0 & 0.015 & 0.068 & -- \\
IM 0.0, repaired and rewritten & data + history & local (thesis) & 0.0 & 100.0 & 0.003 & 0.006 & 0.0 \\
IM 0.0 rewritten + silent exit & data + history & local (thesis) & 0.0 & 99.8 & 0.013 & 0.005 & 2.7 \\
HM 0.7, discovered & data + history & local (thesis) & 0.0 & 99.9 & 0.000 & 0.010 & 1.7 \\
HM 0.7, rewritten & data + history & local (thesis) & 0.0 & 100.0 & 0.001 & 0.009 & 0.0 \\
\bottomrule
\end{tabular}
\caption{Control-flow replica of ProSiT's routing (token game, no time, no resources): every training visit is simulated with its own manifest, three seeds, measures against the training log; the sequence TVD of a row varies by at most 0.006 between the seeds. Routing weights: P(transition fires $\mid$ enabled), estimated by counting on the aligned training log; ``data + history'' uses the planned stops, the flow type and the counts of the activities done so far, as ProSiT's trees do; ``marking'' adds the state of the net. Guard: local = the simple guard on visible stops and Gate Out; look-ahead = also masks silent steps after which the plan cannot be completed; rejection = simulate without guard and keep only visits with the planned number of stops. Rewritten: the same language without silent transitions.}
\label{tab:app_control_flow_replica}
\end{table}


\paragraph{The weights do not see the marking.} On the discovered IM 0.0, adding
the marking to the state of the weights removes every visit without a stop
(15.7\,\% $\rightarrow$ 0) and lowers the sequence TVD from 0.232 to 0.004. Each
skip of a parallel branch is learned on its own, without knowing whether another
branch has already produced a stop. In a rewritten net, every transition has one
input place, so its weight depends on the state anyway, and adding the marking
changes nothing.

\paragraph{Silent transitions alone are not the cause.} The discovered HM 0.7 has
34 silent transitions but no parallel branches, and it generates the same visits
as its rewritten version. The silent steps only make the simple guard leak: on the
discovered HM 0.7 it has to fall back 1.7 times per 1,000 visits.

\paragraph{A guard has to act where the decision is taken.} On the discovered
IM 0.0, a look-ahead guard or rejection sampling gives every truck its planned
number of stops, but the stops are combined wrongly (sequence TVD 0.167 and 0.068,
against 0.006 for the repaired and rewritten net). For a net with parallel
branches, no guard repairs the combination of stops; for a net without them, such
as HM 0.7, the look-ahead guard is enough.

\paragraph{The flower needs the history in the weights.} With weights from
frequencies only, the flower and the discovered ILP 1.00, which has the same
language, generate 48\,\% visits without a stop. With the planned stops and the
history of the visit in the weights, the flower is as good as the other nets
(sequence TVD 0.005). In ProSiT it is not (sequence TVD 0.077). ProSiT's routing
trees, with at most three splits, do not recover the history of a visit as
completely as a count over the full history.

\paragraph{Limits.} The counting estimator is more flexible than ProSiT's trees,
so the replica shows what the representation allows, not what ProSiT learns.
Without time, overlapping stops are not measured. The alignment gives only one
order for steps in parallel branches.
